%------------------------------------------------------------------------------%
% FUNDAMENTAL CONCEPTS IN GENERAL TOPOLOGY                                     %
%------------------------------------------------------------------------------%
% File  : General_Topology.tex                                                 %
% Author: Klaus Hoeflinger, Beethovenstrasse 11, D-73117 Wangen                %
% Email : khoeflinger@t-online.de                                              %
%------------------------------------------------------------------------------%

%------------------------------------------------------------------------------%
% DOCUMENT CLASS and INCLUDE FILES                                             %
%------------------------------------------------------------------------------%
\documentclass[11pt,twoside]{article}
\usepackage{geometry}
\geometry{paperheight=241mm,
          paperwidth=152mm,
          top=22mm,
          bottom=17mm,
          left=20mm,
          right=20mm,
          textwidth=112mm,
          textheight=202mm,
          headheight=10mm,
          headsep=3mm,
          footskip=9mm,
          includehead,
          includefoot,
          marginparsep=0mm,
          marginparwidth=0mm}

\renewcommand{\baselinestretch}{0.945}\normalsize

\AtBeginDocument{\setlength\abovedisplayskip{2mm}}
\AtBeginDocument{\setlength\belowdisplayskip{2mm}}

%------------------------------------------------------------------------------%
% define title formats                                                         %
%------------------------------------------------------------------------------%
\usepackage{titlesec}

\renewcommand{\thesection}{\arabic{section}}
\renewcommand{\sfdefault}{FiraSans}
\renewcommand{\ttdefault}{pcr}
\renewcommand{\rmdefault}{antiqua}

\titleformat{\part}[display]
{\normalfont \large \rmfamily \bfseries \centering}
{\small{\thepart.}}{2pt}{}

\titleformat{\section}[runin]
{\normalfont \rmfamily \bfseries}
{\thesection}{1pt}{.\;}[.]

%\titleformat{\section}
%{\normalfont \bfseries \small \filcenter}
%{\thesection.\;}{0mm}{}

\titleformat{\subsection}[runin]
{\normalfont \itshape}
{}{0mm}{}

\titleformat{\subsubsection}[runin]
{\normalfont \itshape}
{}{}{}{}

\titleformat{\paragraph}[runin]
{\normalfont}
{}{}{}{}

\titleformat{\subparagraph}[runin]
{\normalfont}
{}{}{}{}

\titlespacing*{\part}          { 0pt}{ 0pt}{ 8pt}
\titlespacing*{\section}       { 0pt}{ 7pt}{ 4pt}
\titlespacing*{\subsection}    { 0pt}{14pt}{ 6pt}
\titlespacing*{\subsubsection} { 0pt}{ 6pt}{12pt}
\titlespacing*{\paragraph}     { 0pt}{ 6pt}{12pt}
\titlespacing*{\subparagraph}  { 0pt}{ 0pt}{12pt}

%------------------------------------------------------------------------------%
% define enumerate parameters                                                  %
%------------------------------------------------------------------------------%
\usepackage{enumitem}
\setlength{\parindent}{3pt}
\setlength{\parskip}  {0pt}

%\setlist{noitemsep}

\setlist[enumerate]{
         leftmargin=12pt,
         rightmargin=0pt,
         itemsep=1pt,
         itemindent=3pt,
         topsep=3pt,
         labelsep=4pt,
         partopsep=0pt,
         labelwidth=0pt,
         listparindent=0pt,
         parsep=0pt}

\setlist[itemize]{
         leftmargin=12pt,
         rightmargin=0pt,
         itemsep=0pt,
         itemindent=1pt,
         topsep=1pt,
         labelsep=4pt,
         partopsep=1pt,
         labelwidth=0pt,
         listparindent=0pt,
         parsep=0pt}

%------------------------------------------------------------------------------%
% define packages                                                              %
%------------------------------------------------------------------------------%
\usepackage{upref}
\usepackage{amssymb}
\usepackage{slantsc}
\usepackage{amsmath}
\usepackage{amsthm}
\usepackage{thmtools}
\usepackage{amscd}
\usepackage{verbatim}
\usepackage{latexsym}
\usepackage{multicol}
\usepackage{graphicx}
\usepackage{setspace}
\usepackage[ngerman,english]{babel}
\usepackage[protrusion=true,expansion=true]{microtype}
\usepackage{tikz}
\usetikzlibrary{arrows,arrows.meta,decorations.markings}
\usepackage{mathrsfs}
\usepackage{amsfonts}
\usepackage{datetime2}
\usepackage{textgreek}
\usepackage{hyperref}
\usepackage{pifont}
\usepackage{relsize}
%\usepackage{biblatex}
\relsize{-0.05}
\usepackage[skip=0mm]{parskip}
\usetikzlibrary{decorations.pathmorphing}

\usepackage{mlmodern}
\usepackage[T5,T1]{fontenc}
\usepackage{ragged2e}

%\usepackage{mathabx}
%\usepackage{times}

%\usepackage{fontspec}
%\setmainfont{Nimbus Roman No9 L}
%\usepackage[T1]{fontenc}

%------------------------------------------------------------------------------%
% define page layout                                                           %
%------------------------------------------------------------------------------%
\usepackage{fancyhdr}
\pagestyle{fancy}
\setlength{\headheight}{30.0pt}
\renewcommand{\headrulewidth}{0pt}

\AtBeginDocument{
  \addtolength\abovedisplayskip{-0.0\baselineskip}
  \addtolength\belowdisplayskip{-0.0\baselineskip}
}

\fancyhead{}
\fancyfoot{}
\addtolength{\headwidth}{0.02\marginparwidth}

\makeatletter
\let\ps@plain\ps@fancy
\makeatother

\renewcommand\sectionmark[1]
{
 \def\sectionname{#1}
 \markright{\thesection #1}
}

\makeatletter
\@addtoreset{section}{part}
\makeatother

%------------------------------------------------------------------------------%
% define footnote without number                                               %
%------------------------------------------------------------------------------%
\newcommand{\myfootnote}[1]{%
\renewcommand{\thefootnote}{}%
\footnotetext{#1}%
\renewcommand{\thefootnote}{\arabic{footnote}}%
}

%------------------------------------------------------------------------------%
% define numbering in equations                                                %
%------------------------------------------------------------------------------%
\numberwithin{equation}{section}

%------------------------------------------------------------------------------%
% define newtheorem                                                            %
%------------------------------------------------------------------------------%
\declaretheoremstyle[
headfont=\scshape, 
headindent = 0mm, %\parindent,
%postheadhook = {\hspace{0mm}\newline},
%postheadspace = \newline, 
spaceabove = 3mm, 
spacebelow = 3mm,
numberwithin=section,
name=Theorem]{khMATH}
\declaretheorem[style=khMATH]{theorem}

\declaretheoremstyle[
headfont=\bfseries, 
headindent = 0mm, %\parindent,
%postheadhook = {\hspace{0mm}\newline},
%postheadspace = \newline, 
spaceabove = 3mm, 
spacebelow = 3mm,
numberwithin=part,
name=Theorem]{khMATH1}
\declaretheorem[style=khMATH1]{theorem1}

\declaretheoremstyle[
headfont=\scshape, 
headindent = 0mm, %\parindent,
%postheadhook = {\hspace{0mm}\newline},
%postheadspace = \newline, 
spaceabove = 3mm, 
spacebelow = 3mm,
numberwithin=section,
name=Corollary]{khMATH1}
\declaretheorem[style=khMATH1]{corollary}

%            name         counter     label              numeration-mode
\theoremstyle{plain}
\newtheorem*{introduction}            {\sc Introduction}
\newtheorem {standard}                {}                 [section]
\newtheorem*{definition}              {\sc Definition}
\newtheorem{satz1}                    {\sc Satz}
\newtheorem{lemma}                    {\sc Lemma}
\newtheorem*{proposition}             {\sc Proposition}
%\newtheorem{corollary}               {\sc Corollary}
\newtheorem*{corollar}                {\sc Korollar}
\newtheorem*{remark}                  {\sc Remark}
\newtheorem*{remarks}                 {\sc Remarks}
\newtheorem*{example}                 {Example}
\newtheorem*{examples}                {Examples}
\newtheorem*{theo}                    {\sc Theorem}

%------------------------------------------------------------------------------%
% define tabular                                                               %
%------------------------------------------------------------------------------%
\usepackage{tabularx}
\newcolumntype{L}[1]{>{\raggedright\arraybackslash}p{#1}}
\newcolumntype{C}[1]{>{\centering\arraybackslash}  p{#1}} 
\newcolumntype{R}[1]{>{\raggedleft\arraybackslash} p{#1}} 

%------------------------------------------------------------------------------%
% define \sh, \zB                                                              %
%------------------------------------------------------------------------------%
\def\dsh{{d.\,h.}}
\def\ie{{i.\,e.}}
\def\zB{z.\,B.}
\renewcommand{\abstractname}{ }

%------------------------------------------------------------------------------%
% change default spacing for cases                                             %
%------------------------------------------------------------------------------%
\makeatletter
\renewcommand*\env@cases[1][1.6]{%
  \let\@ifnextchar\new@ifnextchar
  \left\lbrace
  \def\arraystretch{#1}%
  \array{@{}l@{\quad}l@{}}%
}
\makeatother

\newenvironment{rcases}
  {\left.\begin{aligned}}
  {\end{aligned}\right\rbrace}

%------------------------------------------------------------------------------%
% change default for \predisplaypenalty                                        %
%------------------------------------------------------------------------------%
\predisplaypenalty=990
\allowdisplaybreaks \numberwithin{equation}{section}

%------------------------------------------------------------------------------%
% general notations                                                            %
%------------------------------------------------------------------------------%
\def\*{\star}
\def\={\,=\,}
\def\+{\,+\,}
\def\xsubset{{\,\subset\,}}
\def\xtimes{{\,\times\,}}
\def\xeof{{\,\in\,}}
\def\eq{{\,=\,}}
\def\xto{{\,\to\,}}
\def\eqdef{\,:=\,}
\def\:{{\,:\,}}
\def\ele{{\,\in\,}}
\def\exists{\mbox{ exists }}
\def\and{\mbox{ and }}
\def\for{\mbox{ for }}
\def\fa{\mbox{ for all }}
\def\fe{\mbox{ for each }}
\def\qfa{\quad \mbox{ for all }}
\def\df{\,\dotfill}
\def\iff{if and only if }
\def\wrt{with respect to}

\makeatletter
\newcommand \Df {\leavevmode \cleaders \hb@xt@ .70em{\hss .\hss }\hfill \kern \z@}
\makeatother

\def\sql{\mathfrak{a}}               % sesquilinear form a
\def\DF{B}                           % Dirichlet form a
%------------------------------------------------------------------------------%
% number sets and special elements                                             %
%------------------------------------------------------------------------------%
\def\P{\mathbb{H}}                   % half plane of $\C$
\def\J{I}                            % interval I
\def\N{{\mathbb{N}}}                 % unsigned integers
\def\Q{{\mathbb{Q}}}                 % rational integers
\def\Z{{\mathbb{Z}}}                 % integers
\def\R{{\mathbb{R}}}                 % real numbers
\def\Rn{{\mathbb{R}}^n}              % real numbers in Rn
\def\Rd{{\mathbb{R}}^d}              % real numbers in Rn
\def\C{{\mathbb{C}}}                 % complex numbers
\def\F{{\mathbb{K}}}                 % arbitrary field
\def\Torus{{\mathbb{T}}}             % complex circle line
\def\T{\tau}                         % upper bound of interval (0,t)
\def\sign#1{\operatorname{sign}(#1)} % sign of a number

%------------------------------------------------------------------------------%
% set theory                                                                   %
%------------------------------------------------------------------------------%
\def\G{G}                            % domain $G$
\def\clG{\overline{\G}}              % closure of $G$
\def\bndG{\partial \G}               % boundary of $G$
\def\setA{\mathcal{A}}               % set A
\def\setB{\mathcal{B}}               % set B
\def\setC{\mathcal{C}}               % set C
\def\setM{M}                         % set M
\def\setN{N}                         % set N
\def\setS{\mathcal{S}}               % set S
\def\famF{\mathfrak{F}}              % family of sets
\def\indI{\mathcal{I}}               % index set I
\def\pot#1{\mathfrak{P}(#1)}         % power set

\def\relR{\mathbf{R}}                % relation R
\def\mapf{f}                         % mapping f
\def\mapg{g}                         % mapping g

\def\set#1#2{\{ \, #1 \,: \, #2 \, \}}
\def\tensor#1#2{#1 \otimes #2}
\def\Tens{\oplus}

\def\cal#1{{\mathcal{#1}}}
\def\aut#1{\operatorname{Aut}(#1)}

\def\cross#1#2{#1 {\,\times\,} #2}
\def\multicross#1#2{#1 \times  \ldots \times  #2}
\def\multitens#1#2 {#1 \otimes \ldots \otimes #2}

%------------------------------------------------------------------------------%
% group theory                                                                 %
%------------------------------------------------------------------------------%
\def\1{{\mathsf{1}}}                % unit element of a monoid
\def\0{{\mathsf{0}}}                % unit element of an Abelian monoid
\def\kernel#1{\mathfrak{N}(#1)}     % kernel of a homomorphism
\def\cokernel#1{\mathfrak{C}{(#1)}} % cokernel of a homomorphism
\def\Hom#1{\operatorname{Hom}(#1)}

%------------------------------------------------------------------------------%
% linear algebra                                                               %
%------------------------------------------------------------------------------%
\def\vecV{\mathit{V}}               % vector space V
\def\vecH{\mathit{H}}               % vector space H
\def\vecW{\mathit{W}}               % vector space W
\def\vecU{\mathit{U}}               % subspace U
\def\convC{\mathcal{C}}             % convex set C
\def\basH{\mathfrak{B}}             % Hamel base B

\def\LO#1{\mathscr{L}(#1) }         % algebra of linear transformations

\def\scalar#1#2{ \langle #1,#2 \rangle }
\def\trace#1{\operatorname{tr}(#1)}
\def\dim#1{{\operatorname{dim} \, #1 }}
\def\codim#1{{\operatorname{codim} \, #1 }}
\def\dim{\operatorname{dim}}
\def\codim{\operatorname{codim}}
\def\ind{\operatorname{ind}}

\def\genG{\cal{G}}

\def\algA{\cal{A}}
\def\algB{\cal{B}}
\def\algU{\cal{U}}

\def\idI{\cal{I}}

%------------------------------------------------------------------------------%
% topology                                                                     %
%------------------------------------------------------------------------------%
\def\compK{{K}}          % compact set C
\def\openO{\mathcal{O}}  % open set O
\def\openU{\mathcal{U}}  % open set U
\def\U{\mathcal{U}}      % neighbourhood U
\def\V{\mathcal{V}}      % neighbourhood V
\def\d{\mathit{d}}       % metric function d

\def\interior#1{\mathring{#1}}
\def\closure#1{{\overline{#1}}}

\def\topT{\mathcal{T}}   % topology T
\def\topX{X}             % topological space X
\def\topY{Y}             % topological space Y
\def\topZ{Z}             % topological space Z
\def\metrS{M}            % metric space M

\def\grpG{G}             % topological group G
\def\grpA{A}             % topological group A
\def\grpB{B}             % topological group B
\def\grpC{C}             % topological group C
\def\grpH{H}             % topological group H
\def\grpU{U}             % topological group U
\def\grpV{V}             % topological group V
\def\topG{G}             % topological group G
\def\topH{H}             % topological group H
\def\topU{U}             % topological subgroup U

\def\subS{\cal{S}}
\def\riR{R}              % ring R

%------------------------------------------------------------------------------%
% calculus                                                                     %
%------------------------------------------------------------------------------%
\def\supp#1{\operatorname{supp}(#1)}
\def\singsupp#1{\operatorname{sing \, supp}(#1)}

\def\re{\operatorname{Re}}
\def\im{\operatorname{Im}}
\def\grad{\operatorname{grad}}

%\def\abs#1 {\left|#1\right|}
\def\abs#1 {\lvert #1 \rvert}
%\def\abs#1 {\left|\,#1\,\right|}

%------------------------------------------------------------------------------%
% measure theory                                                               %
%------------------------------------------------------------------------------%
\def\salgA{\Sigma}               % sigma-algebra S
\def\salgB{\cal{B}}              % sigma-algebra B
\def\Borel#1{\cal{B}_{#1}}       % Borel-algebra 
\def\LB{\lambda}                 % Lebesgue-Borel measure

%------------------------------------------------------------------------------%
% function spaces                                                              %
%------------------------------------------------------------------------------%
\def\Lp{L^p(\G)}                 % spaces L_p
\def\Lq{L^q(\G)}                 % spaces L_q
\def\Lh{L^2(\G)}                 % spaces L_2
\def\Li{L^{\infty}(\G)}          % spaces L_infty
\def\Lc{L_{loc}^1(\G)}           % spaces L_1^{loc}
\def\Ck{C^k(\G)}                 % spaces C_k

\def\BV#1{BV(#1) }               % set of functions of bounded variation
\def\AC#1{AC(#1) }               % set of absolutely continuous functions

\def\MP#1{\mathscr{M}^+(#1) }    % set of positive measures
\def\MS#1{\mathscr{M}^-(#1) }    % vector space of finite signed measures
\def\MC#1{\mathscr{M}(#1) }      % vector space of complex measures
\def\MCR#1{\mathscr{M}_{r}(#1) } % vector space of regular complex measures

\def\SobW{W^{k,p}(\G)}                % Sobolev space W
\def\SobO{\mathring{W}^{k,p}(\G)}     % Sobolev space W
%\def\WN#1#2{\mathring{W}^{#1}(#2) }   % Sobolev space W
\def\WN#1#2{W_0^{#1}(#2) }   % Sobolev space W
%\def\HN#1#2{\mathring{H}^{#1}(#2) }   % Sobolev space H
\def\HN#1#2{H_0^{#1}(#2) }   % Sobolev space H
%------------------------------------------------------------------------------%
% functional analysis                                                          %
%------------------------------------------------------------------------------%
\def\snorm#1{\lVert #1 \rVert}
\newcommand{\norm}[1]{\left\lVert #1 \right\rVert}

\def\topV{V}                       % topological vector space V
\def\topW{W}                       % topological vector space W
\def\normS{X}                      % normed space

\def\banX{\mathit{X}}              % Banach space X
\def\banY{\mathit{Y}}              % Banach space Y
\def\banZ{\mathit{Z}}              % Banach space Z
\def\dbanX{\mathit{X^*}}           % dual space of Banach space X
\def\dbanY{\mathit{Y^*}}           % dual space of Banach space Y
\def\dbanZ{\mathit{Z^*}}           % dual space of Banach space Z

\def\banA{\mathit{A}}              % Banach algebra A
\def\banB{\mathit{B}}              % Banach algebra B
\def\banC{\mathit{C}}              % C*-algebra C


\def\domain#1{\mathfrak{D}(#1)}    % domain
\def\range#1{\mathfrak{R}(#1)}     % range

\def\isoJ{\mathfrak{J}}            % isometry J
\def\repA{{\cal{A}_{\sql}}}        % representation operator A
\def\linS{{S}}                     % linear operator S
\def\linG{{G}}                     % linear operator G
\def\linL{{L}}                     % linear operator L
\def\linT{{T}}                     % linear operator A
\def\linA{\mathit{A}}              % linear operator A
\def\linB{{B}}                     % linear operator B
\def\linM{{M}}                     % linear operator M
\def\linU{{U}}                     % linear operator U
\def\linI{{I}}                     % linear operator I
\def\A{\cal{A}}                    % linear differential operator A
\def\BDO{\cal{B}}                  % linear boundary differential operator B
\def\linU{{U}}                     % unitary operator U
\def\linP{{P}}                     % projection P
\def\compA{k}                      % compact operator k
\def\linFR{f}                      % finite rank operator f
\def\fredA{A}                      % Fredholm operator F

\def\HAM{\cal{H}}                  % Hamilton operator
\def\PA{\partial^{\alpha}}         % elementary differential operator
\def\DO{\cal{A}}                   % linear differential operator
\def\BC{\cal{B}}                   % boundary operator
\def\SLO{\cal{L}_{p,q}}            % Sturm-Liouville operator
\def\MO{\cal{L}_{M}}               % momentum operator

\def\HS#1{S(#1)           }        % algebra of Hilbert-Schmidt operators
\def\FR#1{F(#1)           }        % algebra of finite rank operators
\def\TC#1{N(#1)           }        % algebra of trace class operators
\def\BU#1{\mathscr{U}(#1) }        % algebra of unitary linear operators
\def\BO#1{\mathscr{L}(#1) }        % algebra of bounded linear operators
\def\CO#1{\mathscr{C}(#1) }        % algebra of compact linear operators
\def\CL#1{\mathscr{C}(#1) }        % algebra of compact linear operators
\def\FO#1{\Phi(#1) }               % set of Fredholm operators
\def\FK#1#2{\Phi_{#2}(#1) }        % set of Fredholm operators with index k

\def\hilH{\mathit{H}}              % Hilbert space H
\def\clU{\mathit{U}}               % closed subspace U

\def\orthO{\mathfrak{O}}           % orthonormal system O
\def\orthS{\mathfrak{S}}           % orthonormal system S
\def\basB{\mathfrak{B}}            % basis B of a Hilbert space

\def\GL#1 {\mathbf{GL}(#1)}        % linear group
\def\OG#1 {\mathbf{O}(#1)}         % linear group
\def\SOG#1 {\mathbf{SO}(#1)}       % linear group

%------------------------------------------------------------------------------%
% distribution theory                                                          %
%------------------------------------------------------------------------------%
\def\disF{F}                  % distribution F
\def\disG{G}                  % distribution G
\def\disH{H}                  % distribution H
\def\disU{U}                  % distribution U
\def\disT{T}                  % distribution T
\def\SF{C^{\infty}(\G)}       % space of smooth functions
\def\TF{C_c^{\infty}(\G)}     % space of compactly supported smooth functions
\def\TFRd{C_c^{\infty}(\R^d)} % space of test functions
\def\SRd{\mathscr{S}(\R^d)}   % Schwartz space
\def\SR1{\mathscr{S}(\R)}     % Schwartz space
\def\B1#1{C^1(#1) }           % space of continuously differentiable functions
\def\Bm#1{C^m(#1) }           % space of continuously differentiable functions
\def\Ck{C^k(\G)}
\def\TD{\mathscr{S}'(\R^d)}   % space of temperated distributions
\def\E{\cal{E}}               % space of distributions with compact support
\def\D{\mathscr{D}}           % D
\def\FT{\mathcal{F}}          % Fourier transformation F
\def\FT{\mathscr{F}}          % Fourier transformation F

%------------------------------------------------------------------------------%
% other definitions                                                            %
%------------------------------------------------------------------------------%
\def\Proof{\textsc{Proof}}
\def\FRECHET{Fr\'{e}chet}
\def\ARZELA{Arzel\'{a}}
\def\GARDING{G\r{a}rding}
\def\HOELDER{H\"older}
\def\SCHROEDINGER{Schr\"odinger}
\def\JOERGENS{J\"orgens}
\def\WUEST{W\"uest}
\def\KAEHLER{K\"aehler}
\def\HOEFLINGER{H\"oflinger}
\def\POINCARE{Poincar\'{e}}
\def\FA{Functional Analysis}

\def\Author   {K.\,{\HOEFLINGER}}
\def\AuthorUC {K.\,HOEFLINGER}
\def\Version  {1.0}
\def\Email    {khoeflinger@t-online.de}


\titleformat{\part}%[display]
{\normalfont \rmfamily \bfseries \centering}
{{\thepart.}}{3pt}{}

\titleformat{\section}[runin]
{\normalfont \bfseries \filcenter}
{\thesection.\,}{0mm}{}[.]

\titlespacing*{\part}   {5pt}{12pt}{ 4pt}
\titlespacing*{\section}{0pt}{ 3pt}{ 3pt}

%------------------------------------------------------------------------------%
%                                BEGIN OF DOCUMENT                             %
%------------------------------------------------------------------------------%
\begin{document}
\thispagestyle{empty}

\centerline{\textrm{\large{\textbf{Fundamental Concepts in General Topology}}}}
\vspace{4pt}
\centerline{\textsc{\small{\Author}}}
\vspace{10pt}

\fancyhead[OL]{\small{\thepage}}
\fancyhead[ER]{\small{\thepage}}

%------------------------------------------------------------------------------%
%                                  CONTENT                                     %
%------------------------------------------------------------------------------%
\myfootnote
{
  Version {\Version} from \today;
  Email:  \textit{\Email}
}

\setcounter{page}{1}

%\newpage
\thispagestyle{empty}

Our text is a short review of the core notations of general topology,
often purely called \textit{point-set topology}.
This mathematical theory deals with the study of limits,
continuity and geometric properties of the underlying sets.
Rather topology has meaning in its own right,
it mainly serves as foundation of further mathematical areas
like \textit{functional analysis}, \textit{differential topology},
\textit{Riemannian geometry} and other ones.
Topological concepts are frequently required to understand
mathematical analysis on a high and abstract level.
One may justly claim
that analysis today is unthinkable
without including the extremely useful notions provided by topology.

\subparagraph{}
The content is divided into four parts: first of all,
we outline the basic concepts of topology.
In the subsequent parts we shall have a look at specific kinds
of topological spaces,
which especially are
\textit{metric spaces} and
\textit{topological groups}.
%and \textit{manifolds}.

\fancyhead[OC]{\small{\textsc{I. Topological Spaces}}}
\fancyhead[EC]{\small{\textsc{I. Topological Spaces}}}
\part{Topological Spaces}
%-------------------------------------------------------------------------------
This first part deals with the main concepts of topology,
and the first one has to consider is that of a \textit{topological space}.

\subparagraph{} %- definition of a topological space
Given a non-empty set $\topX$,
any collection $\topT$ of subsets of $\topX$ is called
a \textit{topology} on $\topX$
iff the following conditions are fulfilled:

\begin{itemize}[leftmargin=10mm,itemsep=1pt]
\item[(T-1)]
both,
the empty set $\emptyset$ and the entire set $\topX$ are elements of $\topT$,

\item[(T-2)]
if $\openO_i \xeof \topT$ for $i \eq 1,\ldots,n$,
then $\openO_1 \cap \ldots \cap \openO_n \xeof \topT$.

\item[(T-3)]
if $\{\openO_{\alpha}\}$ is any collection of members of $\topT$,
then $\bigcup \openO_{\alpha} \xeof \topT$.
\end{itemize}

This means
that a topology is closed under the formation of finitely
many intersections and arbitrarily many unions.
A \textit{topological space} is a pair $(\topX,\topT)$
consisting of a set $\topX$ and a topology $\topT$.
It is convenient to denote the elements of a topological space
as \textit{points}.

%------------------------------------------------------------------------------%
\section{Open and Closed Sets}                                                 %
%------------------------------------------------------------------------------%
For any topological space $(\topX,\topT)$
the members of the topology $\topT$ are called
the \textit{open sets} in $\topX$.
On the other hand,
any subset of $\topX$ is said to be \textit{closed},
if its complementary set is open.
Due to the properties of an open set,
the intersection of arbitrarily many closed sets as well as
the union of a \textit{finite} number of closed sets is closed again.
Especially the base set $\topX$ and the empty set are both closed.

\paragraph{}
The framework of a topological space also enables the definition
of some further important concepts in point-set topology:

\begin{itemize}[leftmargin=15pt,itemsep=1pt]
\item[(a)]
If $\setA$ is a subset of a topological space $\topX$,
then the set $\interior{\setA}$ denotes the union
of all open subsets of $\setA$,
or equivalently,
it is the largest open set $\openO$
such that $\openO \subset \setA$.

\item[(b)]
The set $\interior{\setA}$ is said to be the \textit{interior} of $\setA$,
and each element of $\interior{\setA}$ is called an \textit{inner point}
of $\setA$.

\item[(c)]
Conversely,
consider for a given set $\setA$ the intersection $\overline{\setA}$
of all closed sets containing $\setA$.
Then the set $\overline{\setA}$ is closed again and is called
the \textit{closure} of $\setA$,
and each point $x \xeof \overline{\setA}$ is said to be a \textit{limit point}
of $\setA$.
\end{itemize}

In summary,
for each subset $\setA \subset \topX$ there is always
the sequence of inclusions
\begin{equation}\label{CHAIN}
\emptyset
\subseteq \, \interior{\setA} \,
\subseteq \, \setA            \,
\subseteq \, \overline{\setA}  \,
\subseteq \, \topX.
\end{equation}

We take care of the midth subchain
$\interior{\setA} \, \subseteq \, \setA \, \subseteq \, \overline{\setA}$.
If the inclusions can be replaced by equality,
that is,
$\interior{\setA} \eq \setA$ on the left side,
then and only then the set $\setA$ is open.
Similarly,
if $\setA \eq \overline{\setA}$ on the right side,
then and only then the set $\setA$ is closed.

\subparagraph{} %- discrete sets
Now consider the first inclusion $\emptyset \subseteq \interior{\setA}$
in \eqref{CHAIN}.
If it can be replaced by the equality $\emptyset \eq \interior{\setA}$,
then $\setA$ is called \textit{discrete}.
On the other hand,
in case of $\overline{\setA} \eq \topX$,
the set $\setA$ is called \textit{dense} in $\topX$,
that is,
every point $x \xeof \topX$ can be obtained as a limit point of $\setA$.

\subparagraph{} %- boundary points and boundary
Finally,
for a given subset $\setA \subset \topX$,
the relative complement set
$\partial \setA := \overline{\setA} \setminus \interior{\setA}$
is called the \textit{boundary} of $\setA$,
and each element of $\partial \setA$ is said
to be a \textit{boundary point} of $\setA$.
It can be shown
that a set $\openO \subset \topX$ is open
iff it contains no boundary points,
and a set $\setA \subset \topX$ is closed
iff it contains all its boundary points.
We obtain the boundary of $\setA$ as the intersection
of the sets $\overline{\setA}$ and $\overline{\topX \setminus \setA}$,
which implies that a boundary set is always closed.

%------------------------------------------------------------------------------%
\section{Neighbourhoods}                                                       %
%------------------------------------------------------------------------------%
The notion of \textit{neighbourhood} is one of
the most important one in general topology.
Roughly spoken,
any set constitutes a neighbourhood of a point $x \xeof \topX$,
if it contains all points sufficiently close to $x$.
The exact definition states as follows:
a subset $\U \subset \topX$ is called a \textit{neighbourhood}
of $x \xeof \topX$,
iff $x$ is an interior point of $\U$.
For example,
an open set is a neighbourhood for all of its elements.
The neighbourhoods of a given point $x \xeof \topX$ have the property,
that each finite intersection and each union of arbitrarily many
neighbourhoods of $x$ is also a neighbourhood of $x$.
Moreover,
every superset $\U' \supset \U$ of a neighbourhood of $x$
forms a neighbourhood of $x$.

\subparagraph{} %- convergence of sequences
Neighbourhoods are useful to define the concept of \textit{convergence}
for sequences in a topological space $\topX$.
In particular,
a sequence $(\phi_n)$ is called \textit{convergent} to
an element $x \xeof \topX$,
if for each neighbourhood $\U$ of $x$ there are only
finitely many members of $(\phi_n)$ being outside $\U$.
We will study the convergence of sequences again
in the setting of metric spaces (see Part II).

\subparagraph{} %- neighbourhood bases, first countablility axiom
A \textit{neighbourhood base} of a point $x$ is a family $\famF$
of neighbourhoods of $x$ with the additional property
that each open set $\openO_x$ with $x \xeof \openO_x$
contains a neighbourhood $\U_x \xeof \famF$ as a subset:
$\U_x \subset \openO_x$.
Furthermore,
a topological space is said to satisfy the \textit{first countability axiom},
if each point has a mostly \textit{countable} neighbourhood base.

%------------------------------------------------------------------------------%
\section{Hausdorff's Separation Axiom}                                         %
%------------------------------------------------------------------------------%
Every set $\topX$ containing at least two points
can be employed with different topologies.
Obviously,
the largest one is the \textit{discrete} topology on $\topX$,
which coincides with the power set of $\topX$.
On the other hand,
the coarsest thinkable topology on $\topX$ consists
of the empty set and the set $\topX$ only.

\subparagraph{} %- comparison of various topologies
If $\topT_1,\topT_2$ are topologies on $\topX$,
then $\topT_1$ is called \textit{weaker} than $\topT_2$,
if every open set in $\topT_1$ is also an open set in $\topT_2$,
and at the same time $\topT_2$ is called \textit{stronger} than $\topT_1$.
The notion of topology on a set $\topX$ is extremely general,
and it includes pathological examples of very little interest
such as the topology whose only elements are the sets $\emptyset$ and $\topX$.

\subparagraph{} %- Hausdorff spaces
For most applications in analysis,
it is desirable to impose additional conditions to the topology
that guarantee the existence of sufficiently many open sets.
The most important condition of this kind is
the \textit{Hausdorff separation axiom},
which can be stated as follows:
a topological space $\topX$ is called \textit{Hausdorff space},
if for any points $x,y \xeof \topX$
there are open sets $\openO_x$ and $\openO_y$
such that
$x \xeof \openO_x$, $y \xeof \openO_y$ and
$\openO_x \cap \openO_y \eq \emptyset$.
In other words,
in a Hausdorff space
every pair of points can be separated by suitable open sets
with empty intersection.
Besides many other separation axioms in topology
the Hausdorff one is the most important of them,
because nearly all important topological spaces considered in analysis
are Hausdorff spaces.
The Hausdorff separation axiom ensures two desirable properties
of a topological space:
(1) each singleton $\{x\}$ is closed, and
(2) the limit point of any convergent sequence is unique.

%------------------------------------------------------------------------------%
\section{Continuous Mappings}                                                  %
%------------------------------------------------------------------------------%
The notion of \textit{continuity} of a real-valued function
at a certain point $t \xeof \R$ is well known from elementary calculus.
Using neighbourhoods,
it is possible to extend this notion to arbitrary mappings
between topological spaces $\topX, \topY$ as follows:
a mapping $f \: \topX \xto \topY$ is said to be \textit{continuous at a point}
$x \xeof \topX$,
if for each neighbourhood $\V$ of $f(x)$
there is a neighbourhood $\U$ of $x$ with $f(\U) \subseteq \V$.
Furthermore,
the mapping $f$ is called \textit{continuous},
iff $f$ is continuous at \textit{every} point $x \xeof \topX$.

\subparagraph{} %- another definition of continuous mapping
By means of open sets,
an equivalent definition of continuity is the following one,
which is the most common definition in general topology:
a mapping $f \: \topX \xto \topY$ is called \textit{continuous},
if the preimages $f^{-1}(\openO)$ of open sets $\openO \subseteq \topY$
are also open sets in $\topX$.
Obviously,
the property of a mapping being continuous
depends on the choice of the topology in both spaces.

\subparagraph{} %- homeomorphisms
If a continuous mapping $f \: \topX \xto \topY$ is bijective,
we cannot conclude
that the inverse mapping $f^{-1} \: \topY \xto \topX$ is continuous, too.
This fact justifies the notion of \textit{homeomorphism}
between topological spaces:
a bijective continuous mapping $f \: \topX \xto \topY$
is called a \textit{homeomorphism},
if the inverse mapping $f^{-1}$ is continuous as well.
Thus a homeomorphism carries over exactly the open sets in open sets again.
If there is a homeomorphism between the topological spaces
$\topX$ and $\topY$,
then the spaces are called \textit{homeomorphic} to each other
(denoted $\topX \cong \topY$).
In this case we can talk about structural identity of the spaces $\topX$
and $\topY$ {\wrt} their topological properties.
A property of a topological space is called
a \textit{topological invariant},
if it is preserved by homeomorphic mappings.

%------------------------------------------------------------------------------%
\section{Connected Sets} %- connected sets
%------------------------------------------------------------------------------%
Open sets are suitable to define the notion of \textit{connectedness}
for subsets of a topological space:
a set $C \subseteq \topX$ is said to be \textit{connected},
if there are \textit{not} two open sets $\openO_1,\openO_2$ with
$\openO_1 \cap \openO_2=\emptyset$ and $C \subseteq \openO_1 \sqcup \openO_2$,
otherwise it is called \textit{disconnected}.
Each open and connected set is usually called a \textit{domain}.

\subparagraph{}
Clearly,
the property of a set $C \subset \topX$ being connected
depends on the topology of the underlying space.
Connected sets mapped by a continuous function do not lose this property.
More precisely,
if $\setA \subset \topX$ is connected and the mapping $f \: \topX \xto \topY$
is continuous,
then $f(\setA) \subseteq \topY$ is connected as well.
Connectedness is actually a topological invariant.

\subparagraph{}
The most well-known criterion for connectedness is given
by the following theorem:
\textit{a topological space $\topX$ is connected,
iff $\topX$ and the empty set $\emptyset$ are the only sets,
which are open and closed together}.

\subparagraph{} %- connected components
For a topological space $\topX$,
a \textit{connected component} of a point $x \xeof \topX$ is defined as
the union of all connected sets which contain this point.
Each such set is obviously connected
and cannot be extended to a larger connected set by adding further points.
Since the closure of a connected set is connected again,
the connected components of a topological space are always closed.
They form a disjoint decomposition of the space into connected subsets,
being the coarsest qualitative information about the global structure
of a space.

%------------------------------------------------------------------------------%
\section{Generators of New Topologies}                                         %
%------------------------------------------------------------------------------%
It can be shown
that for a base set $\topX$ the intersection of arbitrarily
many topologies on $\topX$ forms a topology on $\topX$ again.
Thus for a given collection $\setS$ of subsets of $\topX$,
we define the topology $[\setS]$ to be the intersection
of all topologies on $\topX$ containing $\setS$.
The set $\setS$ is denoted a \textit{generator} of the topology $[\setS]$,
and reversely $[\setS]$ is said to be \textit{generated} by $\setS$.

\subparagraph{} %- standard topology in the real line
As a first application of this method we define the \textit{standard topology}
$\topT_{\R}$ on the real line $\R$ as the topology generated
by all half-bounded open intervals,
that is,
$\topT_{\R}$ is the smallest topology on $\R$
containing the half-bounded intervals $(-\infty,a)$ and $(b,\infty)$.

\subparagraph{} %- initial topology
Let $\topX$ be any set and
let $f_i \: \topX \xto \topY_i$ ($i \xeof \indI$) be a family of mappings
into arbitrarily many topological spaces $\topY_i$.
We consider the collection $\setS$ of preimages $f_i^{-1}(\openO)$
for all the open sets $\openO \subset \topY_i$.
Then the topology $[\setS]$ generated by $\setS$ is the coarsest
topology on $\topX$ such that all mappings $f_i$ remain still continuous.
It is called the \textit{initial topology} {\wrt} the family
$\{f_i\}_{i \in \indI}$.

\paragraph{} %- derived topologies
Using initial topologies,
we introduce the \textit{subset topology} and the \textit{product topology}:

\begin{itemize}[leftmargin=12pt,itemsep=2pt]
\item[1.]
Let $\topX$ be a topological space
and $\setM \subset \topX$ be a proper subset of $\topX$.
We define the \textit{subset topology}
$\topT_{\setM}$ on $\setM$
as the initial topology {\wrt} the inclusion mapping
$\setM \hookrightarrow \topX$,
that is,
the subset topology is the smallest one such that
the inclusion of $\setM$ into $\topX$ is continuous.
The pair $(\setM,\topT_{\setM})$ is a topological space in its own right,
and a subset $\openO \subset \setM$ is open {\wrt} $\topT_{\setM}$
iff there is an open set $\openO'$ in $\topX$
such that $\openO \eq \setM \cap \openO'$.

\subparagraph{}
It should be mentioned
that the subset topology in $\setM$ can also be received by definition
of its open sets,
\[
\topT_{\setM} \eqdef \{ \setM \cap \openO': \openO' \in {\topT} \}.
\]
A \textit{continuous embedding} of a space $\topX$ into a space $\topY$
is a mapping $\iota \: \topX \xto \topY$
such that the images of open sets in $\topX$ under $\iota$ exactly
form the subset topology of $\iota(\topX)$ considered as a subspace
of $\topY$.

\item[2.]
The \textit{product topology} $\topT_{\topX \times \topY}$
on the cartesian product $\cross{\topX}{\topY}$ is defined to be
the initial topology {\wrt} the projection mappings
$\pi_X \: \cross{\topX}{\topY} \xto \topX$ and
$\pi_Y \: \cross{\topX}{\topY} \xto \topY$.
In other words,
it is the coarsest topology
for which the both mappings are continuous.
Note
that this definition of the product topology can be extended
to the cartesian product of arbitrarily many topological spaces.
Another way to define the product topology on the set $\cross{\topX}{\topY}$
is given by the following one:
$\topT_{\topX \times \topY}$ is the topology generated by the set of all pairs
$(\openO_1,\openO_2)$ with $\openO_1 \xeof \topX$ and $\openO_2 \xeof \topY$.
\end{itemize}

\subparagraph{} %- standard topology in the set of complex numbers
We can use the product topology to define
the \textit{standard topology} $\topT_{\C}$
on the field $\C$ of complex numbers,
since $\C$ can be considered as the cartesian product $\R \times \R$.
Here,
for each factor of $\R \times \R$ the standard topology
of the space $\R$ is used.

%------------------------------------------------------------------------------%
\section{Compact and Locally Compact Spaces}                                   %
%------------------------------------------------------------------------------%
One of the most useful concepts coming from topology
is that of \textit{compactness}.
It is a property of finiteness defined for subsets of a topological space.
Preliminary,
given a topological space $\topX$,
an \textit{open covering} of a set $\compK \subseteq \topX$
is a collection $\famF$ of open sets such that
\[
\compK \, \subseteq \, \bigcup_{\openO \in {\famF}} \openO.
\]
The set $\compK$ is called \textit{compact},
if \textit{every open covering of $\compK$ has a finite subcovering}.

\subparagraph{}
The meaning of compactness is,
that many proofs are much easier and more transparent,
if the set under consideration is assumed to be compact.
In general,
when examining a topological problem
it is a benefite to consider primarily the compact case if possible.
Note that compactness is an {inner property} of a subset,
that is,
a subset $\compK \subset \topX$ is compact
{\wrt} the topology $\topT$ on $\topX$
\iff it is compact in its subset topology $\topT_{\compK}$.

\subparagraph{} %- relatively compact sets
The property of compactness of a set is inherited to its subsets
supposed these are closed,
that is,
every closed subset of a compact set or space is compact,
too.
On the other hand,
if a topological space $\topX$ is Hausdorff,
then every compact subset of $\topX$ is closed.
We shall call a subset \textit{relatively compact},
if its closure is compact.

\subparagraph{} %- compactness criterion for homeomorphisms
Compactness provides also a helpful criterion to decide
whether a given map $f \: \topX \xto \topY$ is homeomorphic.
In particular,
let $\topX$ be a compact space and $\topY$ be a Hausdorff space,
then each continuous and bijective mapping is a homeomorphism.

\paragraph{} %- locally compactness
Since compactness is a rather ambitious property,
topological spaces having this property are very sparse.
But there is a weaker property,
called \textit{locally compactness},
which is fulfilled by a larger number of topological spaces.
Locally compactness is a property
which has to be satisfied by all the elements of a space:
a topological space $\topX$ is called \textit{locally compact},
if each neighbourhood of any point $x \xeof \topX$
contains a compact neighbourhood.
In case of $\topX$ being a Hausdorff space, $\topX$ is locally compact if
and only if each point $x \xeof \topX$ has a compact neighbourhood.

\subparagraph{} %- locally compact Hausdorff spaces
The class of \textit{locally compact spaces} plays an outstanding role
for investigations in other mathematical areas like
\textit{measure theory} and \textit{abstract harmonic analysis}.
Nearly all locally compact spaces of practical interest
are Hausdorff spaces as well,
so the notion \textit{LCH space} is a popular abbreviation for
a topological space being both,
namely \textit{locally compact} and \textit{Hausdorff}.

%------------------------------------------------------------------------------%
\section{Compact-Open Topologies}                                              %
%------------------------------------------------------------------------------%
Given the two sets $\topX$ and $\topY$,
we denote $\topY^{\topX}$ the set of all mappings $f \: \topX \xto \topY$.
In additon,
if both sets are topological spaces,
let $C(\topX,\topY) \subset \topY^{\topX}$ be the subset
of \textit{continuous} mappings $f \: \topX \xto \topY$.
It is our goal to equip the set $C(\topX,\topY)$ with a suitable topology.

\subparagraph{} %- definition of compact-open topologies
Preliminary,
for a compact subset $\compK \subset \topX$ and an open subset
$\topU \subset \topY$ we introduce the family of mappings
\[
S(\compK,\openO) \eqdef \{f \in C(\topX,\topY): f(\compK) \subset \openO \}.
\]
Then the set $C(\topX,\topY)$ can be furnished with the topology $\topT_{co}$,
which is generated by the collection of all such sets $S(\compK,\openO)$,
called the \textit{compact-open topology} on $C(\topX,\topY)$.
It turns out
that the set $C(\topX,\topY)$ equipped with this topology is a Hausdorff space
whenever $\topY$ is Hausdorff.

\paragraph{}
The establishment of the compact-open topology on sets of continuous mappings
can be seen as a standard method in general topology.
Its meaning is based on the fact that
functions in $C(\topX,\topY)$ are closed together,
if their images on compact subsets are so.

\subparagraph{}
In case of $\topY=\R$ or $\topY=\C$,
let both sets be endowed with their standard topology.
Then for the spaces $C(\topX,\R)$ and $C(\topX,\C)$
we get the compact-open topology,
respectively.
Since $\R$ and $\C$ are Hausdorff spaces according to the standard topologies,
the functions sets $C(\topX,\R)$ and $C(\topX,\C)$ prove to be Hausdorff spaces,
too.

%------------------------------------------------------------------------------%
\section{Baire Spaces}                                                         %
%------------------------------------------------------------------------------%
A topological space is called a \textit{Baire space}
(due to \textsc{R.L.\,Baire}),
if the union of a countable collection of closed sets with empty interior
has empty interior, too.
This definition is equivalent to the following one:
if the union of countable many closed subsets has an interior point,
then at least one of them must have an interior point.
\textsc{Baire}\textit{'s category theorem} claims
(as a special variant of a more general theorem)
that \textit{every LCH space is a Baire space}.
The word "category" is due to \textsc{Baire}'s original terminology,
according to which a set is said to be \textit{of first category}
if it is a countable union of nowhere dense sets,
and \textit{of second category} otherwise.
We may reformule \textsc{Baire}'s theorem by quoting
that \textit{all LCH spaces are of second category}.

%------------------------------------------------------------------------------%
%                                 BIBLIOGRAPHY                                 %
%------------------------------------------------------------------------------%
\begin{comment}
\vspace{8mm}
\begin{small}
\begin{thebibliography}{99}

\bibitem{Du_1} J.\,Dugundji:
\textit{Topology}.
Allyn and Bacon, 1966.

\end{thebibliography}
\end{small}
\end{comment}

%\newpage
%\thispagestyle{empty}
\fancyhead[OC]{\small{\textsc{II. Metric Spaces}}}
\fancyhead[EC]{\small{\textsc{II. Metric Spaces}}}
\part{Metric Spaces}
%-------------------------------------------------------------------------------
The theory of \textit{metric spaces} may serve as a link
between general topology and modern analysis.
Introduced by \textsc{M.\,{\FRECHET}},
they are one of the most important classes within that of topological spaces.
Roughly spoken,
they can be considered as a straightforward generalization
of the Euclidean spaces.
The main property of a metric space is
that a distance value, called the \textit{metric} on that space,
is assigned to each pair of its elements.
This leads to an intuitive understanding
of the topological concept of \textit{closeness} between points.

%------------------------------------------------------------------------------%
\section{Metrics}                                                              %
%------------------------------------------------------------------------------%
Given a non-empty set $\metrS$,
a mapping $\d \: \metrS \xtimes \metrS \xto \R$ is called a \textit{metric}
on $\metrS$
iff the following four conditions hold for all $x,y,z \xeof \metrS$:

\begin{itemize}[leftmargin=17mm,itemsep=1pt]
\item[(M-1)]
$\d(x,y) \geq 0$,

\item[(M-2)]
$\d(x,y) \eq 0$ if and only if $x \eq y$,

\item[(M-3)]
$\d(x,y) \eq \d(y,x)$;

\item[(M-4)]
$\d(x,z) \leq \d(x,y) + \d(y,z)$ \hspace{3mm}(\textit{triangle inequality}).
\end{itemize}

\paragraph{} %- metric spaces
If $\d$ is a metric on the set $\metrS$,
then the pair $(\metrS,\d)$ is called a \textit{metric space}.
Any set $\metrS$ can be made into a metric space in different ways
by employing different metric functions.

\subparagraph{}
There is a natural way to furnish a metric space with a distinguished topology.
Let $(\metrS,\d)$ be a metric space,
then any set of the form
$
B_{\epsilon}(x) := \{y \xeof \metrS: \d(x,y) < \epsilon \}
$
is called an \textit{$\epsilon$-ball} around a point $x \xeof \metrS$.
We then introduce open sets in $\metrS$ by means of $\epsilon$-balls:
a subset $\openO \subseteq \metrS$ is called open,
if for every point $x \xeof \openO$
there is a ball $B_{\epsilon}(x)$
such that $B_{\epsilon}(x) \subset \openO$.

\subparagraph{}
Indeed,
the system $\topT_d$ of all such open sets in $\metrS$ forms a topology
and is called the \textit{induced} or \textit{metric topology} on $\metrS$.
Thus each metric space is naturally a topological space,
and all the concepts provided by general topology can be applied.
Note that metric topologies are always Hausdorff,
since any two points in $\metrS$ can be separated
by sufficient small $\epsilon$-balls.

\paragraph{} %- boundedness
It is easy to define the notion of \textit{boundedness}
for arbitrary subsets of a metric space $(\metrS,\d)$.
For a set $\setA \subseteq \metrS$,
the nonnegative number
\[
\d(\setA) \eqdef \sup \, \{ \d(x,y) \: x,y \in \setA \}
\]
is called the \textit{diameter} of $\setA$,
and $\setA$ is said to be \textit{bounded}
iff $\d(\setA) < \infty$.

\subparagraph{}
Boundedness is \textit{not} a topological invariant of the metric topology.
For example,
consider the real line with the first metric $\d_1(x,y) := \abs{x-y} $
and the second metric
\[
\d_2(x,y) \eqdef \frac{ \abs{x-y} }{1 + \abs{x-y} }
\quad \quad (x,y \in \R).
\]
Then the identity mapping $\iota \:(\R,\d_2) \xto (\R,\d_1)$ is a homeomorphism,
but the diameter of the entire set $\R$ related to the metric $\d_2$ is finite,
whereas {\wrt} the metric $\d_1$ it is infinite.

\paragraph{} %- isometries
Let $(\metrS_1,\d_1)$ and $(\metrS_2,\d_2)$ be metric spaces.
Then a mapping $\iota \: \metrS_1 \xto \metrS_2$ is called an \textit{isometry},
if it satisfies the equations
\[
\d_2(\iota(x),\iota(y)) \= \d_1(x,y)
\quad \quad (x,y \in \metrS_1).
\]

\subparagraph{}
This definition implies that an isometry is injective.
Two metric spaces $\metrS_1,\metrS_2$ are said to be
\textit{isometric} to each other,
if there is a \textit{surjective} isometry between them.
In particular,
isometric spaces are homeomorphic to each other {\wrt} their metric topologies.

%------------------------------------------------------------------------------%
\section{Convergence and Completeness}                                         %
%------------------------------------------------------------------------------%
The concept of convergence is already known
for sequences in arbitrary topological spaces.
By the validness of the first countability axiom
in a metric space $(\metrS,\d)$,
a sequence $(x_n)$ of elements in $\metrS$ converges to
a limit point $x \xeof \metrS$ \iff $\lim_{n \xto \infty} \d(x,x_n)=0$,
that is,
if for each $\epsilon > 0$ there is an integer $N_{\epsilon}$
such that $\d(x_n, x) < \epsilon$ for all $n > N_{\epsilon}$.
Because each metric space is Hausdorff,
the limit point of a convergent sequence must be unique.
By the triangle inequality
\[
\d(x_n,x_m) \, \leq \, \d(x_n,x) + \d(x,x_m),
\]
we can see that a convergent sequence in $\metrS$ satisfies the
\textit{Cauchy convergence condition}
\[
\lim_{n,m \to \infty} \d(x_n,x_m) \= 0.
\]
This condition suggests the following definition:
a sequence $(x_n)$ is called \textit{Cauchy sequence},
if for each $\epsilon > 0$ there is a positive integer $N_{\epsilon}$
such that $\d(x_m, x_n) < \epsilon$ for all $m,n > N_{\epsilon}$.
In addition,
two metrics $\d_1,\d_2$ on a set $\metrS$ are called \textit{equivalent},
if they create identical sets of Cauchy sequences,
and in this case their corresponding metric topologies coincide.

\paragraph{}
It is one of the benefits when working with metric spaces
that the property of a set to be closed or dense can be expressed
in terms of convergent sequences:

\begin{itemize}[leftmargin=12pt,itemsep=1pt]
\item[1.]
a subset $\setA \subset \metrS$ is closed,
iff the limit point of any convergent sequence in $\setA$ belongs to $\setA$,
too,

\item[2.]
the set $\setA$ is dense in $\metrS$,
iff for each point $x \xeof \metrS$ there is a sequence $(x_n)$
in $\setA$ which is convergent to $x$.
\end{itemize}

\subparagraph{}
Clearly,
every convergent sequence is a Cauchy sequence,
but the converse statement does not hold:
consider Cauchy sequences in the space $\Q$ of rational numbers,
which in general do not have a limit point in $\Q$.
This counterexample gives rise to introduce the concept
of \textit{completeness} for a metric space,
which is of outstanding meaning in higher analysis:
indeed,
a metric space $(\metrS,\d)$ is called \textit{complete},
if each Cauchy sequence $(x_n)$ is convergent to an element $x \xeof \metrS$.

\subparagraph{}
One can easily prove
that the cartesian product of arbitrarily many metric spaces is complete
iff each single space is complete,
and that a subset $\A$ of a complete metric space is complete
\iff $\A$ is closed.

\paragraph{}
In the following we shall point out the fundamental assertion
that for each \textit{uncomplete} metric space $(\metrS,\d)$
there is a \textit{complete} metric space $(\overline{\metrS},\delta)$ and
a continuous embedding $\iota \: \metrS \xto \overline{\metrS}$,
which preserves the metric in the sense that
\[
\d(x,y) \= \delta(\iota(x),\iota(y))
\quad \, \quad (x,y \in \metrS).
\]
The metric space $\overline{\metrS}$ is called the \textit{completion}
of $\metrS$ and is unique up to isometric mappings.
Moreover,
the image of $\metrS$ is a dense subset in $\overline{\metrS}$.

\subparagraph{}
The space $\overline{\metrS}$ is usually designed as
a set of equivalence classes consisting of Cauchy sequences in $\metrS$.
More precisely,
two Cauchy sequences $(x_n)$ and $(y_n)$ are called equivalent,
if $\lim_{n \to \infty} \d(x_n,y_n) \eq 0$,
and the mapping
\[
\delta([(x_n)],[(y_n)]) \eqdef \lim_{n \to \infty} \d(x_n,y_n)
\]
defines a metric function on $\overline{\metrS}$,
which is independent of the representants in $[(x_n)]$ and $[(y_n)]$.
The embedding of $\metrS$ into the metric space $\overline{\metrS}$
is performed by mapping an element $x \xeof \metrS$
to that equivalent class $[x]$,
which contains the \textit{constant} Cauchy sequence $(x,x,\ldots)$.
The construction of the real numbers as the completion of the space
of rational numbers is a prototypical example for that method.

\subparagraph{}
Complete metric spaces turn out to be Baire spaces.

%------------------------------------------------------------------------------%
\section{Compactness in Metric Spaces}                                         %
%------------------------------------------------------------------------------%
In the setting of metric spaces 
the property of compactness and relatively compactness of a given subset
can be expressed in terms of \textit{sequentially compactness}
(\textsc{Bolzano-Weierstrass} \textit{theorem}):

\begin{theorem}\label{BOLZANO_WEIERSTRASS_THEOREM}
A subset $\setA \subseteq \metrS$ of a metric space is relatively compact
iff each sequence in $\setA$ has a convergent subsequence.
\end{theorem}

\subparagraph{}
For a compact subset $\compK$ in $\metrS$,
the limit point of such a subsequence must be a member of $\compK$ again,
since $\compK$ is closed.
Furthermore,
every compact metric space is \textit{complete},
due to the fact that
each Cauchy sequence in a compact metric space
must have a convergent subsequence.
In order to describe compact subsets in a metric space,
the following concept is adequate:
a subset $\setA \subseteq \metrS$ is called \textit{totally bounded},
if for each $\epsilon > 0$ there is a finite number of points
$x_1,\ldots,x_n \xeof \setA$
such that their $\epsilon$-balls form an open cover of $\setA$.

\paragraph{}
For \textit{complete} metric spaces the \textit{Heine-Borel} \textit{theorem}
asserts
that there is a strong link between the concepts of totally boundedness
and relatively compactness:

\begin{theorem}
If a metric space $\metrS$ is complete,
then a subset $\setA \subset \metrS$ is relatively compact,
if and only if it is totally bounded.
\end{theorem}

\subparagraph{}
For the Euclidean spaces $\Rd$ the notion of totally boundedness
and usual boundedness are equivalent, 
so a subset $\compK \subset \R^n$ is compact
iff $\compK$ is closed and bounded.

\paragraph{}
The following property of a compact subset in a metric space is often used:
if a set $\setA_1$ is \textit{compact} and
a further set $\setA_2$ is \textit{closed},
then they have non-zero distance to each other,
that is,
$\delta(\setA_1,\setA_2) \eq \delta(\setA_2,\setA_1) > 0$ holds.

%------------------------------------------------------------------------------%
\section{Continuity in the Setting of Metric Spaces}                           %
%------------------------------------------------------------------------------%
In metric spaces
there is an alternative approach to define continuity of a mapping:
a function $f \: D \subseteq \metrS_1 \xto \metrS_2$
is called \textit{continuous} at a point $x_0 \xeof \metrS_1$,
if for each $\epsilon > 0$ there exists $\delta > 0$
such that $d_1(x_0,x) < \delta$ implies $d_2(f(x_0),f(x)) < \epsilon$
for all $x \xeof D$.

\subparagraph{}
Beyond that,
the function $f$ is called \textit{continuous},
if $f$ is continuous at every point of its domain $D$.
Notice
that in metric spaces a function $f$ is continuous at a point $x$,
iff it is \textit{sequentially continuous} in $x$, that is,
a sequence $(x_n)$ converging to $x$ implies the
convergence of the sequence $(f(x_n))$ to the point $f(x)$.
The property of a function being continuous at a point $x$
is of local nature.

\paragraph{} %- Arzela-Ascoli theorem
Let $C(\topX,\C)$ be the set of continuous complex-valued functions
$f \: \topX \xto \C$.
If $\topX$ is a compact topological space,
then any function $f \xeof C(\topX,\C)$ is bounded.
As a consequence,
the difference $f{-}g$ of such functions is bounded,
too,
and one may introduce a metric $\d$ on the set $C(\topX,\C)$ by
\[
\d(f,g) \eqdef \sup_{x \in \topX} \abs{ f(x)-g(x) }
\quad \mbox{for } f,g \in C(\topX,\C).
\]

\subparagraph{}
It is an important task to identify the compact subsets
of the set $C(\topX,\C)$ {\wrt} its metric topology.
More precisely,
the question arises,
under which conditions a subset of $C(\topX)$ is relatively compact.
The main condition in this problem is called \textit{equicontinuity}:
a family $\famF \subset C(\topX)$  is called \textit{equicontinuous},
if for each $\epsilon > 0$ there exists a $\delta > 0$ such that
\[
\abs{(f(x)-f(y)} < \epsilon, \, \mbox{ whenever } \, \d(x,y) < \delta
\]
for every function $f \xeof \famF$.
In addition,
the family $\famF$ is called \textit{uniformly bounded} in $x \xeof \topX$,
if there is a number $M$ such that $\abs{f(x)} \leq M$ for all $f \xeof \famF$.
Once these concepts are available,
the answer to the question above is given
by the \textit{{\ARZELA}-Ascoli theorem}:

\begin{theorem}\label{ARZELA_ASCOLI_THEOREM}
Let $\topX$ be a compact space and $C(\topX)$ the
space of continuous complex-valued functions.
Then a subset $\famF \subset C(\topX)$ is relatively compact,
iff $\famF$ is equicontinuous and uniformly bounded for each $x \xeof \topX$.
\end{theorem}

\subparagraph{}
There are some generalizations of this theorem.
For example,
it is still valid for families of continuous mappings $f \: \topX \xto \metrS$,
where $\topX$ is a compact Hausdorff space and $\metrS$ is a metric space.
Note that in this case the notion of relatively compactness means
that for each point $x \xeof \topX$
the set $\famF_x := \{ f(x): f \xeof \famF \}$ is relatively compact.

\paragraph{} %- uniform continuity
A further notion in metric space theory is \textit{uniform continuity},
which avoids the dependency of the number $\delta$ from the point $x$:
a function $f \: D \subseteq M_1 \xto M_2$
is called \textit{uniformly continuous},
if for each $\epsilon > 0$ there exists $\delta > 0$,
such that $d_1(x,y) < \delta$ implies $d_2(f(x),f(y)) < \epsilon$
for \textit{all} $x,y \xeof D$.
Thus the value $\delta$ can be chosen independently of the points $x,y$.

\paragraph{} %- Heine-Cantor theorem
The \textsc{Heine-Cantor} \textit{theorem} quotes:

\begin{theorem}\label{HEINE_CANTOR_THEOREM}
If $M_1,M_2$ are metric spaces and $M_1$ is \textit{compact},
then every continuous function $f \: M_1 \xto M_2$ is uniformly continuous.
\end{theorem}

\subparagraph{}
Uniformly continuous functions map Cauchy sequences
to Cauchy sequences again,
but this is not true in case of non-uniformly continuous functions.
Furthermore,
the uniform continuity of a function $f$ is a global property,
and each such function clearly is continuous.
The converse statement however is not true:
consider the mapping $f \: \R \xto \R$, $x \mapsto x^2$,
which is continuous,
but \textit{not} uniformly continuous on the real line.

%------------------------------------------------------------------------------%
\section{{\HOELDER}-Continuity}                                                %
%------------------------------------------------------------------------------%
A function $f \: \metrS_1 \xto \metrS_2$ between metric spaces $M_1,M_2$
is called \textit{{\HOELDER}-continuous},
if there a real constants $c \geq 0$ and $\alpha \geq 0$ such that
for all $x,y \xeof \metrS_1$ the so called \textit{{\HOELDER} condition}
\[
\d_2(f(x),f(y)) \, \leq \, c \cdot \d_1(x,y)^{\alpha}
\]
is satisfied.
Moreover,
each {\HOELDER} continuous function $f \: M_1 \xto M_2$
is actually uniformly continuous.
In case of $\alpha \eq 0$,
{\HOELDER} continuity means that the function $f$ is simply bounded,
whereas in case of $\alpha \eq 1$ it is equivalent to the inequality
\[
\d_2(f(x),f(y)) \, \leq \, c \cdot \d_1(x,y),
\]
whose validness is known as \textit{Lipschitz continuity}.
The number $c$ is called the \textit{Lipschitz constant} of the function $f$.

\subparagraph{}
If the Lipschitz constant satisfies $0 < c < 1$,
the function is called a \textit{contraction}.
Contractions in the context of complete metric spaces
are very useful,
because they ensure the existence and uniqueness
of the solution for a couple of nonlinear equations.
This fact is based on the \textsc{Banach} \textit{fixpoint theorem}:

\begin{theorem}\label{BANACH_FIXPOINT_THEOREM}
Let $(\metrS,\d)$ be a complete metric space
and the mapping $f \: \metrS \xto \metrS$ be a contraction.
Then $f$ admits exactly one fixpoint of the equation $f(x) \eq y$.
\end{theorem}

\subparagraph{}
Indeed,
beginning with any $x_0 \xeof \metrS$,
the sequence $(x_n)$ defined by $x_n := f(x_{n-1})$ is Cauchy,
whose limit point $x$ is the solution of the equation $f(x) \eq x$.

%------------------------------------------------------------------------------%
%                                 BIBLIOGRAPHY                                 %
%------------------------------------------------------------------------------%
\begin{comment}
\vspace{8mm}
\begin{small}
\begin{thebibliography}{99}

\bibitem{Du_2} J.\,Dugundji:
\textit{Topology}.
Allyn and Bacon, 1966.

\end{thebibliography}
\end{small}
\end{comment}

%\newpage
%\thispagestyle{empty}
\fancyhead[OC]{\small{\textsc{III. Topological Groups}}}
\fancyhead[EC]{\small{\textsc{III. Topological Groups}}}
\part{Topological Groups}
%-------------------------------------------------------------------------------
Besides the metric spaces there is a further subclass of topological spaces,
which is of great importance in analysis,
namely that of \textit{topological groups}.

\subparagraph{}
First recall from group theory
that an arbitrary group is a set $\topG$ together with
an associative operation $\circ \: \topG \times \topG \xto \topG$, such that

\begin{itemize}[leftmargin=12pt,itemsep=1pt]
\item[1.]
there is exactly one \textit{identity element} $\1 \xeof \topG$ such that
\[
x \circ \1 = \1 \circ x \= x \quad \quad (x \in \topG),
\]

\item[2.]
to each $x \xeof \topG$ there is a unique element $x^{-1} \xeof \topG$,
called the \textit{inverse} of $x$
such that
\[
x \circ x^{-1} = x^{-1} \circ x = \1.
\]
\end{itemize}

\subparagraph{}
If the group operation is commutative,
then the identity element is often denoted $\0$
and each inverse element $x^{-1}$ is denoted $-x$.

\section{Generalities on Topological Groups}
%------------------------------------------------------------------------------%
These are mathematical objects having the algebraic structure of a group
and a topological structure together.
Both kinds must be compatible to each other in the following sense:
let $(\topG,\circ)$ be an ordinary group and
a topological space $(\topG,\topT)$ at once.
Then $\topG$ is called a \textit{topological group},
if the group operation
\[
\circ \: \cross{\topG}{\topG} \xto \topG,
\quad
(x,y) \mapsto x \circ y^{-1}
\]
is a continuous mapping.
Note that
instead to require continuity for the mapping $(x,y) \mapsto x \circ y^{-1}$,
one can equivalently require
that the mappings $(x,y) \mapsto x \circ y$ and $x \mapsto x^{-1}$
be continuous.

\paragraph{} %- examples of topological groups
Well-known examples of topological groups are:

\begin{itemize}[leftmargin=9mm,itemsep=1pt]
\item[1.]
each group $\grpG$ with its discrete topology $2^{\grpG}$,

\item[2.]
the group $\R^{\times} := \R - \{0\}$,

\item[3.]
the group $\C^{\times} := \C - \{(0,0)\}$,

\item[4.]
the group $\R^+ := \set{x \xeof \R}{x >0}$,

\item[5.]
the additive groups $(\R,+)$ and $(\C,+)$,

\item[6.]
each subgroup $\topU$ of a topological group $\topG$
endowed with the subspace topology
and its closure $\overline{\topU}$ are
topological groups in its own right.
\end{itemize}

\paragraph{} %- homomorphisms between topological groups
Like for purely algebraic groups,
a \textit{homomorphism} between topological groups
is defined to be a \textit{continuous} group homomorphism,
and a homomorphism is called an \textit{isomorphism},
if it is continuous and a group isomorphism at the same time.
A well-known example of a group isomorphism is
the \textit{exponential function} $t \mapsto e^{t}$,
which maps the topological group $(\R,+)$ to the group $(\R^+,\cdot)$.
Similarly,
the function
\[
t \in [0,1) \mapsto e^{2 \pi it}
\]
forms a group isomorphism between the groups $\R / \Z$ and $\mathbb{T}$.

\paragraph{}
Since in a topological group the
left  translation $x \mapsto g \circ x$ and the
right translation $x \mapsto x \circ g$
are continuous mappings,
they are both examples for \textit{automorphisms},
that is,
isomorphisms on the group $\topG$ to itself.
Hence the topological properties of a topological group $\topG$
are completely determined by the topology in a neighbourhood
of the identity element $\1 \xeof \topG$,
which gives rise to call a topological group to be \textit{homogeneous}.

\subparagraph{}
The kind,
how the identity $\1$ of a topological group $\topG$
lies in a given subgroup $\topU$,
determines the topology of that subgroup.
In particular,

\begin{itemize}[leftmargin=9mm,itemsep=1pt]
\item[1.]
$\topU \subset \topG$ is open, iff $\1$ is an inner point of $\topU$.

\item[2.]
$\topU$ is discrete, iff $\1$ is an isolated point of $\topU$.
\end{itemize}

\paragraph{} %- circle group
The \textit{circle group} $\mathbb{T}$ is that group
consisting of all complex numbers of modulus $1$ with group operation
\[
e^{i \alpha} \circ e^{i \beta} \eqdef
e^{i \alpha} \cdot e^{i \beta} \eqdef
e^{i(\alpha + \beta)},
\quad \alpha,\beta \in \R.
\]
This group,
endowed with the subset topology inherited
by the standard topology on the field $\C$,
is a topological group.
Consider the set $C(\topG,\mathbb{T})$ consisting of all
continuous functions $f \: \topG \xto \mathbb{T}$.
We equip this set with the compact-open topology,
so $C(\topG,\mathbb{T})$ becomes into a Hausdorff space.

\section{Characters and Dual Groups}
%------------------------------------------------------------------------------%
A \textit{character} is a group homomorphism from a
topological group $\topG$ to the circle group $\mathbb{T}$
(in the sense of topological groups) and hence a continuous mapping.
We denote $\topG^*$ the set of all characters $\chi \: \topG \xto \mathbb{T}$,
which turns out to be a subset of the function space $C(\topG,\mathbb{T})$.
The set $\topG^*$ can be made into a topological space
by choosing the topology $\topT_{\topG^*}$ on $\topG^*$ to be
the subset topology {\wrt} the compact-open topology on $C(\topG,\mathbb{T})$.

\paragraph{} %- definition of the dual group
Now,
performing multiplication of characters
$\chi_1,\chi_2 \xeof \topG^*$ by
\[
\chi_1 \circ \chi_2 \: g \mapsto \chi_1(g) \cdot \chi_2(g),
\]
we obtain a further character $\chi_1 \circ \chi_2 \xeof \topG^*$.
This operation has all the properties of a group operation,
and since it is commutative,
the set $\topG^*$ of all characters of $\topG$
has the structure of an Abelian group,
called the \textit{character group} or the \textit{dual group} of $\topG$.

\subparagraph{}
Furthermore,
one can prove
that the group operation in $\topG^*$ is continuous {\wrt}
the topology $\topT_{\topG^*}$,
hence the dual group of any topological group is a topological group as well.

\paragraph*{} %- examples of classical topological Abelian groups
In the following table we list three classical topological Abelian groups
together with their dual groups and characters:
\vspace{3mm}

\renewcommand{\arraystretch}{1.2}
\centerline
{
\begin{tabular}{|c|c|c|}
\hline
$\topG$ & $\topG^*$ & characters \\
\hline
$\R$ & ${\R}^* \cong \R$ & $\chi_x(t) := e^{itx}, \, x \in \R$ \\
$\Z$ & ${\Z}^* \cong \mathbb{T}$ & $\chi_z(n) := e^{inz}, z \in \mathbb{T}$ \\
$\mathbb{T}$ & ${\mathbb{T}}^* \cong \Z$ & $\chi_n(z) := z^n, \, n \in \Z$ \\
\hline
\end{tabular}
}
\vspace{3mm}
\centerline{\small Table 1: classical topological groups and dual groups.}

\paragraph{}
Concerning cartesian products of topological groups,
for each topological group $\topG$ and $\topH$
we have
\[
(\topG \times \topH)^* \, \cong \, {\topG}^* \times {\topH}^*.
\]
As a consequence,
the group isomorphisms $(\Rd)^* \cong \Rd$ hold for any integer $d$.

\section{Pontryagin Duality}
%------------------------------------------------------------------------------%
The abbreviation \textit{LCA group} usually stands
for a \textit{locally compact Abelian group}.
The meaning of locally compact groups is based on the fact
that they own non-trivial characters,
that is,
we have $\topG^* \neq \{\chi_1\}$ for such a group $\topG$,
where $\chi_1$ is the character with $\chi_1(x) \eq 1$ for all $x \xeof \topG$.

\paragraph{} %- Pontryagin duality theorem
Suppose that the group ${\topG}$ is \textit{locally compact}.
Then it can be shown (for example by the \textsc{{\ARZELA}-Ascoli} theorem)
that the dual group $\topG^*$ is also locally compact
{\wrt} its compact-open topology.
Let $\topG^{**}$ denote the \textit{bidual group}
of a topological group $\topG$,
which is defined to be the dual group of its character group.

\subparagraph{}
One of the main results in the theory of LCA groups is the celebrated
\textit{Pontryagin duality theorem}:

\begin{theorem}\label{PONTRYAGIN_DUALITY_THEOREM}
If $\topG$ is a LCA group,
then $\topG$ and its bidual group $\topG^{**}$ are isomorphic to each other.
\end{theorem}

\subparagraph{}
Notice that
each group element $g \xeof \topG$ gives rise to a character $\chi_x$
on the dual group $\topG^*$ defined by $\chi_x(\xi) := \xi(x)$
for all $\xi \xeof \topG^*$.
Thus the mapping $x \mapsto \chi_x \xeof \topG^{**}$ is an injective group
homomorphism,
and Pontryagin duality means
that the canonical embedding $\iota \: \topG \xto {\topG}^{**}$,
$x \mapsto \chi_x$ is also surjective.

\paragraph{}
It is easy to see
that the groups listed in table 1 are examples for Pontryagin duality.
This kind of duality also implies remarkable relationships
between the topological properties of the LCA group $\topG$
and its dual group $\topG^*$:

\begin{itemize}[leftmargin=9mm,itemsep=1pt]
\item[1.]
$\topG$ is compact  if and only if ${\topG}^*$ is discrete.

\item[2.]
$\topG$ is discrete if and only if ${\topG}^*$ is compact.
\end{itemize}

\begin{comment}
\section{Homogene Räume}
%------------------------------------------------------------------------------%
Sei im folgenden $(\topG,T,\topX)$ ein G-Raum,
$p \in \topX$ ein fest gewählter Punkt,
$[p] \in \topX/\topG$ der Orbit von $p$ unter $T$
und $\grpG_p$ seine Isotropiegruppe.
Es ist leicht zu zeigen,
dass die Zuordnung $\topG/\topG_p \mapsto [p]$ stetig und bijektiv ist.

\paragraph{}
Im Falle, dass $\topG$ kompakt und $\topX$ ein Hausdorff-Raum ist,
ist auch die Umkehrabbildung $\phi^{-1}$ stetig,
und für jeden Punkt $p \in \topX$ ist der Orbit $[p]$
zum Raum $\topG/\topG_p$ homöomorph.
Die Orbits können also unter diesen Voraussetzungen
als Quotientenräume $\topG/\topG_p$ angesehen werden.

\paragraph{}
Einfach gesprochen heißt eine topologischer Raum {\em homogen},
falls je zwei Punkte durch einen Homöomorphismus ineinander
übergeführt werden können.
Mit anderen Worten,
der topologische Raum ist der Orbit einer transitiven Gruppenoperation und
homöo\-morph zum Quotienten einer topologischen Gruppe
nach einer geeigneten Untergruppe.
Wir definieren daher den Begriff des {\em homogenen Raumes} wie folgt:

\paragraph{}
Sei $\topG$ eine topologische Gruppe, $\topU \subset \topG$ eine Untergruppe
und $\topG/\topU$ die Menge der Linksnebenklassen,
versehen mit der Quotienten-Topologie.
Dann heißt der topologische Raum $\topG/\topU$ ein {\em homogener Raum}.

\begin{itemize}[leftmargin=15pt]
\item[1.]
Der Torus $\T := \R / \Z$ ist ein homogener Raum,
der homöomorph zum Einheitskreis $S^1 \subset \C$ ist.

\item[2.]
Die Menge $S^1 \subset \C$ ist eine kompakte topologische Gruppe,
die auf dem Hausdorff-Raum $\C$ durch folgende Operation stetig operiert:
$T(e^{i \phi},z) := e^{i \phi} \cdot z$.
Die Isotropie-Gruppe für ein Element $z \in \C^{\times}$ ist
die ein-elementige Menge $\enumeration{1} \subset S^1$,
also gilt $\topG/\topG_z \cong S^1$.
Die Räume $[z]$ und $\topG/\topG_z$ sind zueinander homöomorph,
und folglich stellt die Menge $S^1$ einen homogenen Raum dar.

\item[3.]
Jeder Quotientenraum $\R^{k+m} / \R^k$ heißt
eine {\em Graß\-mann-Mannig\-faltig\-keit}.
Die Graß\-mann-Mannig\-faltig\-keit ist also die Menge
aller $k$-dimensionaler Untervektorräume von $\R^{k+m}$.
\end{itemize}

\paragraph{}
Die topologischen Eigenschaften der homogenen Räume $\topG/\topU$ hängen
auf folgende Weise mit denen der Untergruppen $\topU$ab:

\begin{theorem}
Sei $\topG$ eine topologische Gruppe und $U$ eine Untergruppe von $\topG$.
Dann gilt:

\begin{itemize}[leftmargin=15pt]
\item[1.]
$\topG/\topU$ ist genau dann ein Hausdorff-Raum,
wenn $\topU$ abgeschlossen in $\topG$ ist.
\item[2.]
Die Topologie von $\topG/\topU$ ist genau dann diskret,
wenn $\topU$ offen in $\topG$ ist.
\end{itemize}
\end{theorem}

\paragraph{}
Dieser Satz kommt insbesondere dann zur Anwendung,
wenn die Untergruppe $\topU$ ein Normalteiler von $\topG$ ist.
Sei $\overline{\topU}$ der Abschluss von $\topU$,
dann ist der homogene Raum $\topG/\overline{\topU}$ eine
Hausdorffsche topologische Gruppe,
und die Projektionsabbildung $\pi:\topG \to \topG/\overline{\topU}$ ist
ein Epimorphismus topologischer Gruppen.
\end{comment}

%------------------------------------------------------------------------------%
%                                 BIBLIOGRAPHY                                 %
%------------------------------------------------------------------------------%
\begin{comment}
\vspace{8mm}
\begin{small}
\begin{thebibliography}{99}

\bibitem{Ru_3} W.\,Rudin:
\textit{Fourier Analysis on Groups}.
Wiley-Interscience, 1962.

\end{thebibliography}
\end{small}
\end{comment}

%\newpage
%\thispagestyle{empty}
%\fancyhead[OC]{\small{\textsc{IV. Manifolds}}}
%\fancyhead[EC]{\small{\textsc{IV. Manifolds}}}
%\part{Manifolds}
%-------------------------------------------------------------------------------

%------------------------------------------------------------------------------%
%                                 BIBLIOGRAPHY                                 %
%------------------------------------------------------------------------------%
\vspace{8mm}
\begin{thebibliography}{99}

\bibitem{Du} J.\,Dugundji:
\textit{Topology}.
Allyn and Bacon, 1966.

\bibitem{Ru} W.\,Rudin:
\textit{Fourier Analysis on Groups}.
Wiley-Interscience, 1962.

\end{thebibliography}

%------------------------------------------------------------------------------%
%                               END OF DOCUMENT                                %
%------------------------------------------------------------------------------%
\end{document}
